\section{From prime witnesses to every natural number}
\label{sec:assembly}

We now avoid induction on the quantified statement ``$n$ has a
four-square representation.'' A second explicit bounded loop carries
the tuple that is being constructed.

For $n>0$, start with
\[
 q_0=n,\qquad A_0=(1,0,0,0).
\]
For a state $(q,A)$ define
\begin{equation}
\label{eq:factorstep}
 T(q,A)=
 \begin{cases}
 (q,A),&q\le1,\\
 (q/p,\;A*P(p)),&q>1,
 \end{cases}
\end{equation}
where $p$ in the second branch is the first divisor of $q$ in
$2,\ldots,q$. This is the bounded search of Lemma~\ref{lem:lpf}.
Iterate $T$ exactly $n$ times.

\begin{proposition}[Factor-loop invariant]
\label{prop:factor}
For each $n>0$, quantifier-free induction proves
\begin{equation}
\label{eq:factorinvariant}
 1\le q_t,\qquad q_t\Nrm{A_t}=n,\qquad
 (q_t\monus1)\le(n\monus1)\monus t.
\end{equation}
Consequently $q_n=1$ and $\Nrm{A_n}=n$.
\end{proposition}

\begin{proof}
The initial norm is $1$, so the invariant holds at $t=0$.
If $q=1$, the state is unchanged. If $q>1$, Lemma~\ref{lem:lpf}
supplies a prime $p$ dividing $q$, with $2\le p\le q$.
The exact quotient $q'=q/p$ is positive and satisfies $q'<q$:
from $q=pq'$ and $p\ge2$ one obtains $q\ge2q'>q'$.

Theorem~\ref{thm:prime} and Euler's identity give
\[
 q'\Nrm{A*P(p)}
 =q'\Nrm{A}\Nrm{P(p)}
 =q'\Nrm{A}p
 =q\Nrm{A}=n.
\]
Thus the first two parts of the invariant are preserved, and the rank
$q\monus1$ decreases in the saturated sense of Lemma~\ref{lem:clock}.
That lemma supplies the third part. At $t=n$ its bound is zero,
so positivity gives $q_n=1$, after which the norm equation gives
$\Nrm{A_n}=n$.
\end{proof}

\begin{proof}[Proof of Theorem~\ref{thm:main}]
For $n=0$ define $W_i(0)=0$. For $n>0$ take the tuple $A_n$ produced
by~\eqref{eq:factorstep} and define
\[
 W_i(n)=\abs{(A_n)_i}\qquad(0\le i<4).
\]
All components are primitive recursive, by composition of the bounded
searches, coded tuple operations, and the two fixed-clock recursions.
Proposition~\ref{prop:factor} proves that their squares sum to $n$.
This also holds at zero by definition. The split $n=0$ or $n>0$ is
a decidable arithmetic split, giving~\eqref{eq:main} in $\ipra$.

For each coordinate, $W_i(n)^2\le n$ because the other summands are
nonnegative. If $W_i(n)=0$ its bound is immediate. Otherwise write
$W_i(n)=u+1$; then $(u+1)^2\ge u+1$, so $W_i(n)\le n$.
Finally introduce the four existential quantifiers using the four
specific terms $W_i(n)$. This derives~\eqref{eq:exists} intuitionistically.
\end{proof}

\subsection{A compact program specification}

The mathematical program is the following. Its arguments, records, signed
operations, and parity tests are all primitive recursive after the coding
of Appendix~\ref{app:coding}.

\begin{lstlisting}
prime_rep(p):
    if p = 2: return (1,1,0,0)
    if p < 3: return (0,0,0,0)
    h := (p-1) div 2
    search a,b in 0..h for (a*a+b*b+1) mod p = 0
    if the bounded search fails: return (0,0,0,0)
    x := (a,b,1,0); m := norm(x) div p
    repeat exactly p times:
        if m <= 1: retain (m,x)
        else if m is even:
            x := parity_halve(x); m := m div 2
        else:
            r := centered_residues(x,m)
            z := x * conjugate(r)
            (m,x) := (norm(r) div m, z div m)
    return x

four_squares(n):
    if n = 0: return (0,0,0,0)
    q := n; A := (1,0,0,0)
    repeat exactly n times:
        if q = 1: retain (q,A)
        else:
            p := first divisor of q in 2..q
            (q,A) := (q div p, A * prime_rep(p))
    return componentwise_abs(A)
\end{lstlisting}

Simultaneous assignment in this specification uses the old state on every
right-hand side. The distributed Python implementation skips repeated
terminal states by default; its \code{literal\_clock=True} option executes
all the prescribed iterations. Both variants implement the same function
on natural inputs, because the skipped states are absorbing.

\subsection{A still simpler canonical witness definition}

The constructive proof also justifies a very small, though slower, witness
program. Search the $(n+1)^4$ codes for four digits in base $n+1$ and
return the first $(a,b,c,d)$ with $a^2+b^2+c^2+d^2=n$.
Theorem~\ref{thm:main} supplies a successful candidate, since all four
coordinates are at most $n$. Hence the failure sentinel is impossible.
This gives four primitive recursive witnesses defined using natural
arithmetic alone. The descent construction is what proves the bounded
search always succeeds; merely writing the bounded search would not by
itself prove the theorem.

\section{Why the whole argument belongs to \texorpdfstring{$\ipra$}{iPRA}}
\label{sec:audit}

The arithmetic proof above is already constructive. The remaining
foundational question is the strength of its induction. The following
table records every type of induction that carries substantive work.

\begin{longtable}{@{}L{0.22\textwidth}L{0.29\textwidth}L{0.43\textwidth}@{}}
\toprule
Task & Recursion bound & Quantifier-free invariant or equation\\
\midrule
\endhead
Quotient and remainder & Dividend $a$ & $a=mQ+R$ and $R<m$ for $m>0$.\\
Finite search & Prescribed candidate bound & Stored success is a genuine hit;
all earlier failures are a Boolean fold.\\
B\'ezout computation & $p+1$ & Two coefficient identities, positive first
remainder, decreasing second remainder, finite common-divisor bit.\\
Residue counting & $p+1$ rows, $p$ columns & Equations of finite sums and
prefix counts; range and distinctness use finite Boolean folds.\\
Prime descent & $p$ & $I_p(s_t)$ and
$(m_t\monus1)\le(m_0\monus1)\monus t$.\\
General assembly & Original $n$ & $q_t\ge1$, $q_t\Nrm{A_t}=n$, and
$(q_t\monus1)\le(n\monus1)\monus t$.\\
\bottomrule
\end{longtable}

Each state field is a primitive recursive term in the loop counter and
the input parameters. Consequently the norm equalities and rank
inequalities in the last two rows are literally quantifier-free after
expanding the state projections. The prime condition is a fixed
primitive recursive bit equation. Its presence as a parameter assumption
does not add a quantifier to the induction formula.

The finite counting argument may be displayed in ordinary mathematical
notation using statements about all entries of a list. Its actual
implementation uses bounded sums of equality-test bits. Their laws are
proved by induction on numerical prefix lengths. This is a finite
combinatorial proof over coded data, requiring no comprehension axiom.

All uses of contradiction concern a decidable branch already under
consideration: the failed search flag, a zero new multiplier, or a
forbidden divisor. General classical reductio is unnecessary. All
witnesses used after such a split are either already stored in a record
or obtained by exact arithmetic. There is no application of countable
choice, Markov's principle, unbounded minimization, or induction on an
existential representation predicate.

\subsection{Explicit bounds on intermediate arithmetic}

The bound $m<p$ gives $\Nrm{x}=mp<p^2$, hence $\abs{x_i}<p$ throughout
prime descent. In its odd branch, $\abs{r_i}<p$ and each raw coordinate
of $x*\bar r$ has absolute value less than $4p^2$. In its even branch,
raw sums and differences have absolute value less than $2p$.
After division, the norm invariant again bounds the new coordinates by
$p$. The seed search has at most $((p+1)/2)^2$ candidates for odd $p$.

In the outer loop, $q\Nrm{A}=n$ and $q\ge1$ give
$\Nrm{A}\le n$, so every accumulated coordinate has absolute value at
most $n$. Every selected prime is at most $q\le n$.
Before normalization, a product coordinate is bounded in absolute value
by $4n^2$. A fixed nesting of polynomial-size pair codes remains
primitive recursive. These elementary register bounds are useful for
mechanization, though ordinary primitive recursion does not require
polynomial bounds on its intermediate values.

The construction is intended to expose logical strength, not to optimize
running time in the binary length of $n$. Under a unit-cost convention for
integer arithmetic and division, the direct fixed-clock program has a
crude $O((n+1)^3)$ bound on its arithmetic operations: at most $n$ outer
stages, each with a search of at most $(n+1)^2$ candidates and at most $n$
descent steps. This observation is not a polynomial-time claim in
$\log n$ and is not needed for the proof.
